\section{A sharp congruence bound for nonmaximal image orders}
\label{order:sec:main}

\subsection{Setup and main theorem}

Let $\OO$ be a complete discrete valuation ring with fraction field $F$,
uniformizer $\pi$, and residue field $k$. Normalize $v(\pi)=1$ and set
$v(0)=\infty$. Let $G$ be an arbitrary group and
\[
 \rho,\sigma:G\longrightarrow\GL_d(\OO),\qquad d\ge1.
\]
Assume that their rational matrix spans are full:
$\Span_F\rho(G)=\Mat_d(F)=\Span_F\sigma(G)$.
The integral spans
\[
 A_\rho=\Span_{\OO}\rho(G),\qquad A_\sigma=\Span_{\OO}\sigma(G)
\]
are orders in $\Mat_d(F)$. Define
\[
 \begin{aligned}
 e_\rho&=\min\{j\ge0:\pi^j\Mat_d(\OO)\subseteq A_\rho\},\\
 e_\sigma&=\min\{j\ge0:\pi^j\Mat_d(\OO)\subseteq A_\sigma\},\qquad
 e=\min(e_\rho,e_\sigma).
 \end{aligned}
\]
Here ``conductor exponent'' means this containment exponent; it is not
an Artin conductor. Let
\[
 a=\min_{g\in G}v\bigl(\tr\rho(g)-\tr\sigma(g)\bigr),\qquad
 H_n=\{X\in\Mat_d(\OO/\pi^n):\sigma(g)X=X\rho(g)\ \forall g\in G\}.
\]
We use $\min(n,\infty)=n$. It is enough in all upper bounds below to
replace $e_\rho,e_\sigma$ by certified upper bounds and $a$ by the
valuation of any one nonzero trace discrepancy.

\begin{theorem}[Sharp trace-and-order bound]\label{thm:main}
Under these assumptions, for every integer $n\ge1$,
\begin{equation}\label{eq:main}
 \boxed{\quad \len_{\OO}H_n\le
 \min(n,a)+(d^2-1)\min(n,e).\quad}
\end{equation}
If $a<\infty$, then the rational representations are inequivalent and
\begin{equation}\label{eq:total}
 C(\rho,\sigma):=\sup_n\len_{\OO}H_n
 \le a+(d^2-1)e.
\end{equation}
Write the relation Smith exponents, including zero exponents, as
$0\le b_1\le\cdots\le b_{d^2}$, with an exponent $\infty$ for each free
summand. There is at most one infinite exponent, and
\begin{equation}\label{eq:profile}
 b_{d^2-1}\le e\quad(d\ge2),\qquad
 b_{d^2}\le a+e\quad(a<\infty).
\end{equation}
Consequently $\pi^eH_n$ is cyclic for every $n$. Both the coefficient
$d^2-1$ in \eqref{eq:total} and the loss $e$ in the last inequality of
\eqref{eq:profile} are sharp.
\end{theorem}

The statement extends the article's exact maximal-order theorem rather
than replacing it: when $e=0$ the bound becomes $\len H_n\le\min(n,a)$.
Carayol's trace-rigidity argument gives equality in that case; see also
Serre's exposition \cite[Theorem 1]{Serre1993}. For $e>0$
there need not be equality, and traces alone do not determine the module.
The defect contribution in \eqref{eq:main} stops growing once $n\ge e$.

\begin{corollary}[A scalar discrepancy]\label{cor:scalar-discrepancy}
Suppose in addition that for some $g\in G$,
$\rho(g)=\alpha I_d$ and $\sigma(g)=\beta I_d$, with $\alpha\ne\beta$.
Put $\delta=v(\alpha-\beta)$. Then, independently of possible trace
cancellation when the residue characteristic divides $d$,
\[
 \len H_n\le\min(n,\delta)+(d^2-1)\min(n,\delta,e),\qquad
 C\le\delta+(d^2-1)\min(\delta,e).
\]
These bounds are sharp for every $\delta,e\ge1$ and $d\ge2$.
\end{corollary}
\begin{proof}
The $g$ equation makes every invariant factor at most $\delta$. Combine
this with \eqref{eq:profile}. The family of Example~\ref{ex:length},
with its scalar exponent $a$ replaced by arbitrary $\delta\ge1$ without
requiring $\delta\ge e$, has Smith profile
$\delta,\min(\delta,e),\ldots,\min(\delta,e)$ and attains the bound.
The residue prime need not be prime to $d$ for this scalar statement.
\end{proof}

\subsection{Proof}

\subsubsection{Relation modules and their least exponent}

Put $m=d^2$ and $V=\Mat_d(\OO)$. For every $g\in G$ let
\[
 D_g(X)=\sigma(g)X-X\rho(g).
\]
Let $R\subseteq V^*$ be generated by all coordinate functions of all $D_g$.
Because $V^*$ is Noetherian, finitely many of these functions generate $R$.
Choose a matrix $D$ whose rows are such a finite set. Smith normal form gives
\[
 V^*/R\simeq\bigoplus_{i=1}^{m}\OO/\pi^{b_i},
\]
where a factor with $b_i=0$ is zero and a factor with $b_i=\infty$ denotes
$\OO$. In particular
\begin{equation}\label{eq:length}
 \len H_n=\sum_{i=1}^m\min(n,b_i).
\end{equation}
There is an $\OO$-basis $C_1,\ldots,C_m$ of $V$ such that every $D_g(C_i)$
has all entries in $\pi^{b_i}\OO$ when $b_i<\infty$, and vanishes when
$b_i=\infty$. For finite $b_i$ some relation has valuation exactly $b_i$.
Every basis matrix $C_i$ is primitive: some entry is a unit.

If $s=b_1<\infty$, every entry of every $D_g(X)$ is divisible by $\pi^s$
for every integral $X$. Taking $X=I_d$ gives
\begin{equation}\label{eq:mintrace}
 s\le a.
\end{equation}
If all relations are zero, full rational span forces $d=1$ and
$\rho=\sigma$, and the theorem is immediate. We henceforth exclude this
degenerate case when using $s$.

\subsubsection{The singular-intertwiner obstruction}

\begin{lemma}\label{lem:singular}
Let $X\in\Mat_d(\OO)$ be primitive and singular. If
$D_g(X)\in\pi^r\Mat_d(\OO)$ for all $g$, then $r\le e$.
The same conclusion holds after a finite extension of the valued field,
with valuations still normalized by $v(\pi)=1$.
\end{lemma}
\begin{proof}
Choose a primitive vector $v\in\OO^d$ with $Xv=0$ by clearing denominators
in a nonzero vector in the kernel over $F$. Then
$X\rho(g)v\in\pi^r\OO^d$ for every $g$, hence
$XAv\in\pi^r\OO^d$ for every $A\in A_\rho$. Choose $j$ with $v_j$ a
unit. Since $\pi^{e_\rho}E_{ij}\in A_\rho$,
\[
 \pi^{e_\rho}v_jXe_i\in\pi^r\OO^d\qquad(1\le i\le d).
\]
If $r>e_\rho$, every entry of $X$ is in the maximal ideal, contradicting
primitivity. Applying the same argument to $X^{\mathsf T}$ and the
transposed equations gives $r\le e_\sigma$. These prove $r\le e$.
Over the valuation ring of a finite extension the order containments
remain true after scalar extension, and the same argument applies.
\end{proof}

\begin{lemma}[One direction above the defect]\label{lem:pencil}
For every integer $r>e$, the image of $H_r$ under reduction to
$\Mat_d(k)$ has dimension at most one. In particular
$b_{m-1}\le e$ when $m>1$.
\end{lemma}
\begin{proof}
Suppose two classes give linearly independent residual matrices and lift
them to $X,Y\in\Mat_d(\OO)$. Every linear combination satisfies the
intertwining equations modulo $\pi^r$. Over an algebraic closure of $F$,
the homogeneous polynomial $\det(\lambda X+\mu Y)$ of degree $d>0$
has a projective zero. This is immediate if it is the zero polynomial;
otherwise every positive-degree homogeneous polynomial in two variables
has a zero in the projective line. The zero is defined over a finite
extension. Normalize its coordinates so that $\lambda,\mu$ are integral
and at least one is a unit. Independence of the residual matrices survives
residue-field extension, so $Z=\lambda X+\mu Y$ is primitive. It is
singular and intertwines modulo $\pi^r$, contradicting
Lemma~\ref{lem:singular}. In Smith coordinates the indicated image has
dimension exactly $\#\{i:b_i\ge r\}$. Take $r=e+1$.
\end{proof}

The pencil argument is why the result concerns absolute matrix span.
Over a residue field without extension, a two-dimensional pencil can
avoid singular matrices at all rational projective points. Passing to a
finite extension is essential and is permitted because the order
containment remains valid there.

\subsubsection{The longest direction and the missing unit of loss}

Assume first that $a<\infty$. A nonzero rational intertwiner between
two absolutely irreducible representations of equal dimension is an
isomorphism, which would force equal traces. Thus the rational Hom space
vanishes and all $b_i$ are finite. Write $b=b_m$ and choose the primitive
Smith basis matrix $X=C_m$. If $b\le e$, then \eqref{eq:mintrace} and
\eqref{eq:length} already give
\[
 C=\sum_i b_i\le a+(m-1)e.
\]

Suppose now $b>e$. Lemma~\ref{lem:singular} shows that $X$ is invertible
over $F$. Let
$0=t_1\le\cdots\le t_d=t$ be its matrix Smith exponents; the first is
zero because $X$ is primitive. If $t>0$, use integral invertible row and
column operations to write $X=P\operatorname{diag}(\pi^{t_i})Q$, with
unit factors absorbed into $P,Q$, and set
\[
 Y=P E_{dd}Q,\qquad Z=X-\pi^tY.
\]
The matrix $Z$ is singular and remains primitive: the zero Smith exponent
occurs in a different position because $t>0$. Since every relation on
every integral matrix is divisible by $\pi^s$,
\[
 D_g(Z)\in\pi^{\min(b,t+s)}\Mat_d(\OO)\qquad(g\in G).
\]
Lemma~\ref{lem:singular} and $b>e$ therefore imply
\begin{equation}\label{eq:compensation}
 t+s\le e.
\end{equation}
This also proves $t\le e$. If $t=0$, that inequality is immediate and
the extra compensation is unnecessary.

Now $X^{-1}\in\pi^{-t}\Mat_d(\OO)$, and
\[
 \sigma(g)-X\rho(g)X^{-1}=D_g(X)X^{-1}
 \in\pi^{b-t}\Mat_d(\OO).
\]
Taking traces yields
\begin{equation}\label{eq:trace-long}
 b\le a+t\le a+e.
\end{equation}
If $t=0$, Lemma~\ref{lem:pencil} gives
$C\le b+(m-1)e\le a+(m-1)e$.
If $t>0$, then $m>1$ and
\[
 C\le b+s+(m-2)e
 \le (a+t)+(e-t)+(m-2)e
 =a+(m-1)e.
\]
Thus the order loss that can enlarge the longest invariant factor is
paid back by a smaller shortest factor. Bounding the factors separately
would lose one unnecessary copy of $e$.

For $a=\infty$, Lemma~\ref{lem:pencil} still bounds all but one
exponent by $e$, including the assertion that there is at most one free
factor. It gives $\len H_n\le n+(m-1)\min(n,e)$ directly.

Finally let $a<\infty$ and fix $n$. If $n\le e$, use $s\le a$ to get
$\len H_n\le\min(n,a)+(m-1)n$. If $n>e$ and $a\ge n$, use
Lemma~\ref{lem:pencil}. If $n>e$ and $a<n$, use the total bound just
proved. This establishes \eqref{eq:main} in every case. The cyclicity of
$\pi^eH_n$ follows from the Smith profile. This completes the proof of
the bounds in Theorem~\ref{thm:main}.

\subsection{Sharpness and exact examples}

\subsubsection{The exponent loss is necessary}

\begin{example}\label{ex:exponent}
Let $\OO=\mathbb Z_p$, let $a,e\ge1$, and let $G$ be free on $u,v,z$.
Set
\[
\begin{array}{lll}
 \rho(u)=I+p^eE_{12},&\rho(v)=I+E_{21},&\rho(z)=I,\\
 \sigma(u)=I+E_{12},&\sigma(v)=I+p^eE_{21},&
 \sigma(z)=\operatorname{diag}(1,1+p^a).
\end{array}
\]
Both image-order conductor exponents equal $e$. The $u,v$ equations force
\[
 X=\operatorname{diag}(x,p^ex)
\]
over every quotient ring, with no additional torsion coordinates. The
$z$ equation is $p^{a+e}x=0$. Consequently the relation Smith exponents are
$(0,0,0,a+e)$, and
\[
 \len H_n=\min(n,a+e).
\]
The trace depth is $a$. Indeed, conjugation of $\rho$ by
$\operatorname{diag}(1,p^e)$ agrees with $\sigma$ on $u,v$ and agrees
modulo $p^a$ on $z$. These conjugated matrices are integral, so their
traces agree modulo $p^a$ on every word and its inverse. The trace
discrepancy at $z$ is $-p^a$.

For the order assertions, $A_\rho$ has basis
$I,E_{21},p^eE_{12},p^eE_{11}$. The analogous order for $\sigma$ has basis
$I,E_{12},p^eE_{21},p^{\min(a,e)}E_{11}$. They are closed under
multiplication, contain the generators and their inverses, and are
generated by their indicated differences and products. The off-diagonal
condition proves that their conductor exponent is exactly $e$.
\end{example}

\subsubsection{The length coefficient is necessary}

\begin{example}\label{ex:length}
Let $d\ge2$, choose a prime $p\nmid d$, and take $a\ge e\ge1$. Let $G$
be free on $z$ and $x_{ij}$, $1\le i,j\le d$. Put
\[
 \rho(x_{ij})=\sigma(x_{ij})=I+p^eE_{ij},\quad
 \rho(z)=I,\quad\sigma(z)=(1+p^a)I.
\]
Every generator is invertible over $\mathbb Z_p$. Both image orders are
$\OO I+p^e\Mat_d(\OO)$, with conductor exponent $e$. The centralizer
relations have one zero Smith factor and $d^2-1$ factors $p^e$:
commuting with every matrix unit means precisely that all off-diagonal
entries vanish and all diagonal entries are equal. The scalar $z$
relation therefore gives
\[
 (b_1,\ldots,b_{d^2})=(e,\ldots,e,a),\qquad
 \len H_n=\min(n,a)+(d^2-1)\min(n,e).
\]
The trace discrepancy at $z$ is $-dp^a$, of valuation $a$; all
discrepancies are divisible by $p^a$. This proves equality in the main
bound and the optimality of its coefficient $d^2-1$.
\end{example}

\subsubsection{The complete defect profile for scalar twists}

\begin{proposition}\label{prop:scalar}
Let $\tau:J\to\GL_d(\OO)$ have full rational span, and let the
commutator relation module for $\tau$ have invariant decomposition
\[
 Q_{\rm ad}\simeq\OO\oplus\bigoplus_{i=1}^{d^2-1}\OO/\pi^{c_i},
 \qquad c_i\ge0.
\]
Let $u\in\OO^\times$, $u\ne1$, and put $b=v(u-1)$. On
$G=J\times\mathbb Z$ define
\[
 \rho(h,r)=\tau(h),\qquad\sigma(h,r)=u^r\tau(h).
\]
If $t$ is the valuation of the ideal generated by
$\{\tr\tau(h):h\in J\}$, then the trace depth is $a=b+t$ and
\[
 \len H_n=\min(n,b)+\sum_i\min(n,b,c_i),\qquad
 C=b+\sum_i\min(b,c_i).
\]
In particular this is an exact bad-order replacement for trace depth,
with the full centralizer-defect profile retained.
\end{proposition}
\begin{proof}
The equations for $J$ are commutation with $\tau$, and the equation for
the generator of $\mathbb Z$ is $(u-1)X=0$. Thus $H_n$ is the
$\pi^b$-torsion of the finite-level commutant. Apply its Smith formula.
Every trace discrepancy is $(1-u^r)\tr\tau(h)$; the ideal generated by
$1-u^r$ for integers $r$ is $(1-u)$, since $r=1$ occurs and the geometric
sum works also for negative $r$ using the unit $u$. This proves $a=b+t$.
The trace ideal is nonzero because the matrix span is full and the trace
map on $\Mat_d(F)$ is nonzero in every characteristic.
\end{proof}

\subsection{Finite certificates and arithmetic consequences}

\subsubsection{What a finite computation certifies}

If $G$ is generated by a finite set $S$, the intertwining equations for
$S$ already imply the equations for every group element over every
quotient ring. For products use
\[
 D_{gh}(X)=\sigma(g)D_h(X)+D_g(X)\rho(h),
\]
and for inverses use
$D_{g^{-1}}(X)=-\sigma(g)^{-1}D_g(X)\rho(g)^{-1}$.
Hence a Smith form of the stacked generator equations computes the exact
module, with no enumeration of long words. This requires generators of
the joint image of $(\rho,\sigma)$; generators of just one image are
insufficient.

To certify an upper bound $e_\rho$, choose finitely many words whose
matrices span a lattice $L\subseteq A_\rho$ of full rank in
$\Mat_d(\OO)$. The largest Smith exponent of the coordinate matrix for
$L$ is an upper bound for $e_\rho$. The same applies to $\sigma$.
Any exact trace discrepancy provides an upper bound for $a$. Substituting
these three certified quantities into Theorem~\ref{thm:main} is sound,
even if the finite matrix samples do not generate the whole image order
and even if the trace sample does not attain the minimum.

When entries admit exact arithmetic, a complete image-order basis can be
computed by starting with $\OO I$ and adjoining products by the generator
matrices and their inverses until the module stabilizes. Noetherianity
ensures termination. A stable module is the whole order because it
contains the identity and is stable under all generator words. This is
an algebraic algorithm for an explicitly presented joint image, not a
claim that arbitrary Galois images can be determined effectively.

\subsubsection{The conjugate-dual term of the supplied article}

For the article's pair $\sigma=D(\rho)$, inversion and conjugation
permute group elements, transposition preserves the full matrix lattice,
and multiplying individual group matrices by units leaves their span
unchanged. Thus the two conductor exponents coincide. With that common
value $e$, a nonzero trace discrepancy of valuation $a$ certifies
\[
 C(T)\le a+(d^2-1)e.
\]
More accurately, use the minimum trace depth if it is known. This
quantifies the bad coefficient-prime correction that the article leaves
open. It does not remove the integral matrix-span hypothesis from every
other use in the source: local cocycle ranks, scalar inertia, and the
automorphic twisting construction still need their own arguments.
If the source's scalar element is available with discrepancy depth
$\delta$, Corollary~\ref{cor:scalar-discrepancy} instead gives
$C(T)\le\delta+(d^2-1)\min(\delta,e)$, recovering the full-span
cyclic bound at $e=0$ without tracing a scalar matrix.

For a compatible system over a number field $L$, suppose a fixed
nonzero ideal $J$ is generated by exact common trace discrepancies and
the local orders have certified conductor exponents $e_\lambda$. At all
places where the hypotheses hold,
\[
 C_\lambda\le v_\lambda(J)+(d^2-1)e_\lambda.
\]
If only finitely many $e_\lambda$ are nonzero, put
$\mathfrak f=\prod_\lambda\lambda^{e_\lambda}$, excluding any designated
places as necessary. Then
\[
 \sum_\lambda C_\lambda\log N\lambda
 \le\log N(J)+(d^2-1)\log N(\mathfrak f).
\]
This is exactly the additive form needed for the source's product-formula
budget, with a separately visible correction at bad coefficient primes.

